\documentclass[12pt,aspectratio=169,ignorenonframetext]{ltxbeamer}

\usepackage[spanish,es-nodecimaldot,es-tabla]{babel}

\title{Depuración}
\subtitle{Curso de {\lmr\LaTeX}}
\author{Joar Buitrago\texorpdfstring{\thanks{jebuitragoc@unal.edu.co}}{ <jebuitragoc@unal.edu.co>}}
\date{}

\addbibresource{bibliography.bib}



\begin{document}
\begin{frame}
\maketitle
\end{frame}






\section{Anatomía de un archivo \texttt{.tex}}
\begin{frame}[fragile]
\frametitle{Anatomía de un archivo \texttt{.tex}}

\begin{columns}
    \begin{column}{0.5\textwidth}

        \begin{minted}[escapeinside=||]{latex}
\documentclass[12pt]{article}|\tikz[remember picture,baseline]{\node [anchor=base] (class) {\strut};}|

\title{Mi Primer Documento}|\tikz[remember picture,baseline]{\node [anchor=base] (preamble begin) {\strut};}|
\author{Joar Buitrago}
\date{}|\tikz[remember picture,baseline]{\node [anchor=base] (preamble end) {\strut};}|

\begin{document}
\maketitle|\tikz[remember picture,baseline]{\node [anchor=base] (body begin) {\strut};}|
¡Hola, Mundo {\LaTeX}!|\tikz[remember picture,baseline]{\node [anchor=base] (body end) {\strut};}|
\end{document}
        \end{minted}
    \end{column}

    \begin{column}{0.5\textwidth}
        \tikz[remember picture,baseline]{\coordinate [anchor=base] (braces);}
    \end{column}
\end{columns}

\pause

\begin{tikzpicture}[remember picture, overlay]
    \draw [thick,decoration=brace,decorate] (class.north -| braces) -- (class.south -| braces) node [midway, right] {Clase};
    \draw [thick,decoration=brace,decorate] (preamble begin.north -| braces) -- (preamble end.south -| braces) node [midway, right] {Preámbulo};
    \draw [thick,decoration=brace,decorate] (body begin.north -| braces) -- (body end.south -| braces) node [midway, right] {Cuerpo};
\end{tikzpicture}
\end{frame}





\section{Compilación}
\begin{frame}
\frametitle{Compilación}

\begin{center}
\begin{tikzpicture}
    \begin{scope}[local bounding box=tex]
        \useasboundingbox (-2ex,-2.5ex) rectangle (2ex,2.5ex);
        \draw [very thick]
            [rounded corners=2pt] (-1.5ex,-2ex) coordinate (element) -- ++(0ex,4ex)
            [rounded corners=1pt] -- ++(2ex,0ex) -- ++(1ex,-1ex)
            [rounded corners=2pt] -- ++(0ex,-3ex) -- cycle
            (element) ++(2ex,4ex) |- ++(1ex,-1ex);
        \draw (element) ++(3ex,0ex) node [fill=black, inner sep=0.5ex] {\texttt{TEX}};
    \end{scope}

    \pause

    \begin{scope}[local bounding box=compiler, shift={([xshift=2cm]tex)}]
        \node {\lmr\LaTeX};
    \end{scope}

    \draw [->, very thick] (tex) -- (compiler);

    \pause

    \begin{scope}[local bounding box=pdf, shift={([xshift=2cm]compiler)}]
        \useasboundingbox (-2ex,-2.5ex) rectangle (2ex,2.5ex);
        \draw [very thick]
            [rounded corners=2pt] (-1.5ex,-2ex) coordinate (element) -- ++(0ex,4ex)
            [rounded corners=1pt] -- ++(2ex,0ex) -- ++(1ex,-1ex)
            [rounded corners=2pt] -- ++(0ex,-3ex) -- cycle
            (element) ++(2ex,4ex) |- ++(1ex,-1ex);
        \draw (element) ++(3ex,0ex) node [fill=black, inner sep=0.5ex] {\texttt{PDF}};
    \end{scope}

    \draw [->, very thick] (compiler) -- (pdf);

    \pause

    \begin{scope}[local bounding box=log, shift={([yshift=2cm]pdf)}]
        \useasboundingbox (-2ex,-2.5ex) rectangle (2ex,2.5ex);
        \draw [very thick]
            [rounded corners=2pt] (-1.5ex,-2ex) coordinate (element) -- ++(0ex,4ex)
            [rounded corners=1pt] -- ++(2ex,0ex) -- ++(1ex,-1ex)
            [rounded corners=2pt] -- ++(0ex,-3ex) -- cycle
            (element) ++(2ex,4ex) |- ++(1ex,-1ex);
        \draw (element) ++(3ex,0ex) node [fill=black, inner sep=0.5ex] {\texttt{LOG}};
    \end{scope}

    \begin{scope}[local bounding box=aux, shift={([yshift=-2cm]pdf)}]
        \useasboundingbox (-2ex,-2.5ex) rectangle (2ex,2.5ex);
        \draw [very thick]
            [rounded corners=2pt] (-1.5ex,-2ex) coordinate (element) -- ++(0ex,4ex)
            [rounded corners=1pt] -- ++(2ex,0ex) -- ++(1ex,-1ex)
            [rounded corners=2pt] -- ++(0ex,-3ex) -- cycle
            (element) ++(2ex,4ex) |- ++(1ex,-1ex);
        \draw (element) ++(3ex,0ex) node [fill=black, inner sep=0.5ex] {\texttt{AUX}};
    \end{scope}

    \draw [->, very thick] (compiler) -- (log);
    \draw [->, very thick] (compiler) -- (aux);
\end{tikzpicture}
\end{center}
\end{frame}





\begin{frame}[fragile]
\frametitle{Compilación}

\begin{minted}[escapeinside=\%\%]{text}
%\textit{latex}% %\textit{documento}%.tex
\end{minted}
\end{frame}





\begin{frame}[fragile]
\frametitle{Compilación}

\begin{minted}[escapeinside=\%\%]{text}
lualatex %\textit{documento}%.tex
\end{minted}
\end{frame}





\begin{frame}[fragile]
\frametitle{Compilación}

\begin{minted}[escapeinside=\%\%]{text}
lualatex --interaction=batchmode %\textit{documento}%.tex
\end{minted}
\end{frame}





\begin{frame}[fragile]
\frametitle{Compilación}

\begin{minted}{text}
user@linux$ ls
  documento.aux documento.log documento.pdf documento.tex
\end{minted}
\end{frame}





\begin{frame}[fragile]
\frametitle{Compilación}

\begin{minted}{text}
user@windows> dir
 El volumen de la unidad C no tiene etiqueta.
 El número de serie del volumen es: XXXX-XXXX

 Directorio de C:\Users\User\Documents\Documento

20/08/1803  16:04    <DIR>          .
20/08/1803  16:04    <DIR>          ..
20/08/1803  16:04                32 documento.aux
20/08/1803  16:04             3 434 documento.log
20/08/1803  16:04            31 119 documento.pdf
20/08/1803  16:02               168 documento.tex
               4 archivos         34 753 bytes
               2 dirs  255 485 902 848 bytes libres
\end{minted}
\end{frame}





\section{El archivo \texttt{.log}}
\begin{frame}[fragile]
\frametitle{El archivo \texttt{.log}}

\tiny
\begin{minted}{text}
This is LuaHBTeX, Version 1.24.0 (TeX Live 2026)  (format=lualatex 2026.5.24)  24 MAY 2026 16:04
 restricted system commands enabled.
**documento.tex
(./documento.tex
LaTeX2e <2025-11-01>
L3 programming layer <2026-04-28>
Lua module: luaotfload 2024-12-03 v3.29 Lua based OpenType font support
Lua module: lualibs 2023-07-13 v2.76 ConTeXt Lua standard libraries.
Lua module: lualibs-extended 2023-07-13 v2.76 ConTeXt Lua libraries -- extended collection.
luaotfload | conf : Root cache directory is "/var/lib/overleaf/tmp/texmf-var/luatex-cache/generic/names".
luaotfload | init : Loading fontloader "fontloader-2023-12-28.lua" from kpse-resolved path "/usr/local/texlive/2026/texmf-dist/tex/luatex/luaotfload/fontloader-2023-12-28.lua".
Lua-only attribute luaotfload@noligature = 1
luaotfload | init : Context OpenType loader version 3.134
Inserting `luaotfload.node_processor' in `pre_linebreak_filter'.
Inserting `luaotfload.node_processor' in `hpack_filter'.
Inserting `luaotfload.glyph_stream' in `glyph_stream_provider'.
Inserting `luaotfload.define_font' in `define_font'.
Lua-only attribute luaotfload_color_attribute = 2
luaotfload | conf : Root cache directory is "/var/lib/overleaf/tmp/texmf-var/luatex-cache/generic/names".
Inserting `luaotfload.harf.strip_prefix' in `find_opentype_file'.
Inserting `luaotfload.harf.strip_prefix' in `find_truetype_file'.
Removing  `luaotfload.glyph_stream' from `glyph_stream_provider'.
Inserting `luaotfload.harf.glyphstream' in `glyph_stream_provider'.
Inserting `luaotfload.harf.finalize_vlist' in `post_linebreak_filter'.
Inserting `luaotfload.harf.finalize_hlist' in `hpack_filter'.
Inserting `luaotfload.cleanup_files' in `wrapup_run'.
Inserting `luaotfload.harf.finalize_unicode' in `finish_pdffile'.
Inserting `luaotfload.glyphinfo' in `glyph_info'.
Lua-only attribute luaotfload.letterspace_done = 3
Inserting `luaotfload.aux.set_sscale_dimens' in `luaotfload.patch_font'.
Inserting `luaotfload.aux.set_font_index' in `luaotfload.patch_font'.
Inserting `luaotfload.aux.patch_cambria_domh' in `luaotfload.patch_font'.
Inserting `luaotfload.aux.fixup_fontdata' in `luaotfload.patch_font_unsafe'.
Inserting `luaotfload.aux.set_capheight' in `luaotfload.patch_font'.
Inserting `luaotfload.aux.set_xheight' in `luaotfload.patch_font'.
Inserting `luaotfload.rewrite_fontname' in `luaotfload.patch_font'.
Inserting `tracingstacklevels' in `input_level_string'.
(/usr/local/texlive/2026/texmf-dist/tex/latex/base/article.cls
Document Class: article 2025/01/22 v1.4n Standard LaTeX document class
(/usr/local/texlive/2026/texmf-dist/tex/latex/base/size12.clo
File: size12.clo 2025/01/22 v1.4n Standard LaTeX file (size option)
luaotfload | db : Font names database loaded from /var/lib/overleaf/tmp/texmf-var/luatex-cache/generic/names/luaotfload-names.luc.gz)
\c@part=\count274
\c@section=\count275
\c@subsection=\count276
\c@subsubsection=\count277
\c@paragraph=\count278
\c@subparagraph=\count279
\c@figure=\count280
\c@table=\count281
\abovecaptionskip=\skip49
\belowcaptionskip=\skip50
\bibindent=\dimen149
)
(/usr/local/texlive/2026/texmf-dist/tex/latex/l3backend/l3backend-luatex.def
File: l3backend-luatex.def 2026-02-18 L3 backend support: PDF output (LuaTeX)
\l__color_backend_stack_int=\count282
Inserting `l3color' in `luaotfload.parse_color'.) (./documento.aux)
\openout1 = documento.aux

LaTeX Font Info:    Checking defaults for OML/cmm/m/it on input line 7.
LaTeX Font Info:    ... okay on input line 7.
LaTeX Font Info:    Checking defaults for OMS/cmsy/m/n on input line 7.
LaTeX Font Info:    ... okay on input line 7.
LaTeX Font Info:    Checking defaults for OT1/cmr/m/n on input line 7.
LaTeX Font Info:    ... okay on input line 7.
LaTeX Font Info:    Checking defaults for T1/cmr/m/n on input line 7.
LaTeX Font Info:    ... okay on input line 7.
LaTeX Font Info:    Checking defaults for TS1/cmr/m/n on input line 7.
LaTeX Font Info:    ... okay on input line 7.
LaTeX Font Info:    Checking defaults for TU/lmr/m/n on input line 7.
LaTeX Font Info:    ... okay on input line 7.
LaTeX Font Info:    Checking defaults for OMX/cmex/m/n on input line 7.
LaTeX Font Info:    ... okay on input line 7.
LaTeX Font Info:    Checking defaults for U/cmr/m/n on input line 7.
LaTeX Font Info:    ... okay on input line 7.
LaTeX Font Info:    External font `cmex10' loaded for size
(Font)              <14.4> on input line 8.
LaTeX Font Info:    External font `cmex10' loaded for size
(Font)              <7> on input line 8.
LaTeX Font Info:    External font `cmex10' loaded for size
(Font)              <12> on input line 9.
LaTeX Font Info:    External font `cmex10' loaded for size
(Font)              <8> on input line 9.
LaTeX Font Info:    External font `cmex10' loaded for size
(Font)              <6> on input line 9.
[1

{/usr/local/texlive/2026/texmf-var/fonts/map/pdftex/updmap/pdftex.map}] (./documento.aux)
 ***********
LaTeX2e <2025-11-01>
L3 programming layer <2026-04-28>
 ***********
)

Here is how much of LuaTeX's memory you used:
 462 strings out of 475892
 100000,460012 words of node,token memory allocated
 406 words of node memory still in use:
   3 hlist, 1 vlist, 1 rule, 2 glue, 3 kern, 1 glyph, 4 attribute, 48 glue_spec,
 4 attribute_list, 1 write nodes
   avail lists: 2:80,3:14,4:4,5:28,6:6,7:218,8:1,9:46,10:3,11:4
 23177 multiletter control sequences out of 65536+600000
 37 fonts using 2300247 bytes
 35i,6n,37p,124b,207s stack positions out of 10000i,1000n,20000p,200000b,200000s
</usr/local/texlive/2026/texmf-dist/fonts/opentype/public/lm/lmroman8-regular.otf>
</usr/local/texlive/2026/texmf-dist/fonts/opentype/public/lm/lmroman12-regular.otf>
</usr/local/texlive/2026/texmf-dist/fonts/opentype/public/lm/lmroman17-regular.otf>
Output written on documento.pdf (1 page, 8615 bytes).

PDF statistics: 29 PDF objects out of 1000 (max. 8388607)
 16 compressed objects within 1 object stream
 0 named destinations out of 1000 (max. 1048576)
 1 words of extra memory for PDF output out of 10000 (max. 100000000)
\end{minted}
\end{frame}





\begin{frame}[fragile]
\frametitle{El archivo \texttt{.log}}

\begin{minted}[breaklines]{text}
This is LuaHBTeX, Version 1.24.0 (TeX Live 2026)  (format=lualatex 2026.5.24)  24 MAY 2026 16:04
\end{minted}

\bigskip

Indica qué compilador se está usando, su versión, la versión de los archivos de formato de {\lmr\LaTeX} y la fecha y hora en la que se ejecutó el programa.
\end{frame}





\begin{frame}[fragile]
\frametitle{El archivo \texttt{.log}}

\begin{minted}[breaklines]{text}
 restricted system commands enabled.
\end{minted}

\bigskip

Indica el nivel de acceso que tienen las ejecuciones de programas externos.\pause

Se modifica con la opción \texttt{-shell-escape} o \texttt{-enable-write18}.
\end{frame}





\begin{frame}[fragile]
\frametitle{El archivo \texttt{.log}}

\begin{minted}[breaklines]{text}
**documento.tex
\end{minted}

\bigskip

Indica el archivo que se va a procesar.
\end{frame}





\begin{frame}[fragile]
\frametitle{El archivo \texttt{.log}}

\begin{minted}[breaklines]{text}
(./documento.tex
\end{minted}

\bigskip

Indica que {\lmr\LaTeX} encontró el archivo e inició el procesamiento del mismo.
\end{frame}





\begin{frame}[fragile]
\frametitle{El archivo \texttt{.log}}

\begin{minted}[breaklines]{text}
LaTeX2e <2025-11-01>
\end{minted}

\bigskip

Indica cuándo fue compilado el motor de {\lmr\LaTeXe}.
\end{frame}





\begin{frame}[fragile]
\frametitle{El archivo \texttt{.log}}

\begin{minted}[breaklines]{text}
L3 programming layer <2026-04-28>
\end{minted}

\bigskip

Indica cuándo fue compilado el motor de {\lmr\LaTeX\kern.15em3}.
\end{frame}





\begin{frame}[fragile]
\frametitle{El archivo \texttt{.log}}

\begin{minted}[breaklines]{text}
Lua module: luaotfload 2024-12-03 v3.29 Lua based OpenType font support
Lua module: lualibs 2023-07-13 v2.76 ConTeXt Lua standard libraries.
Lua module: lualibs-extended 2023-07-13 v2.76 ConTeXt Lua libraries -- extended collection.
\end{minted}

\bigskip

Como estamos usando {\lmr Lua\LaTeX}, indica que se han cargado distintos módulos de {\lmr Lua\LaTeX}.
\end{frame}





\begin{frame}[fragile]
\frametitle{El archivo \texttt{.log}}

\begin{minted}[breaklines]{text}
luaotfload | conf : Root cache directory is "/var/lib/overleaf/tmp/texmf-var/luatex-cache/generic/names".

.
.
.

Inserting `tracingstacklevels' in `input_level_string'.
\end{minted}

\bigskip

Como estamos usando {\lmr Lua\LaTeX}, muestran los registros de la carga de los módulos de {\lmr Lua\LaTeX}.
\end{frame}





\begin{frame}[fragile]
\frametitle{El archivo \texttt{.log}}

\begin{minted}[breaklines]{text}
(/usr/local/texlive/2026/texmf-dist/tex/latex/base/article.cls
\end{minted}

\bigskip

Indica que {\lmr\LaTeX} encontró e inició el procesamiento de la clase.
\end{frame}





\begin{frame}[fragile]
\frametitle{El archivo \texttt{.log}}

\begin{minted}[breaklines]{text}
Document Class: article 2025/01/22 v1.4n Standard LaTeX document class
\end{minted}

\bigskip

Informa cuál es la clase que se va a procesar, la fecha de su última modificación, su versión y descripción.
\end{frame}





\begin{frame}[fragile]
\frametitle{El archivo \texttt{.log}}

\begin{minted}[breaklines]{text}
(/usr/local/texlive/2026/texmf-dist/tex/latex/base/size12.clo
\end{minted}

\bigskip

Indica que {\lmr\LaTeX} encontró e inició el procesamiento de una opción de la clase.
\end{frame}





\begin{frame}[fragile]
\frametitle{El archivo \texttt{.log}}

\begin{minted}[breaklines]{text}
File: size12.clo 2025/01/22 v1.4n Standard LaTeX file (size option)
\end{minted}

\bigskip

Informa cuál es la opción que se va a procesar, la fecha de su última modificación, su versión y descripción.
\end{frame}





\begin{frame}[fragile]
\frametitle{El archivo \texttt{.log}}

\begin{minted}[breaklines,breakanywhere]{text}
luaotfload | db : Font names database loaded from /var/lib/overleaf/tmp/texmf-var/luatex-cache/generic/names/luaotfload-names.luc.gz)
\end{minted}

\bigskip

Como estamos usando {\lmr Lua\LaTeX}, muestran los registros de la carga de los módulos de {\lmr Lua\LaTeX}. Además del fin del procesamiento del archivo de opciones.
\end{frame}





\begin{frame}[fragile]
\frametitle{El archivo \texttt{.log}}

\begin{minted}[breaklines]{text}
\c@part=\count274
\c@section=\count275
...
\abovecaptionskip=\skip49
\belowcaptionskip=\skip50
\bibindent=\dimen149
)
\end{minted}

\bigskip

Indica configuraciones iniciales de la clase y el final del procesamiento de la misma.
\end{frame}





\begin{frame}[fragile]
\frametitle{El archivo \texttt{.log}}

\begin{minted}[breaklines,breakanywhere]{text}
(/usr/local/texlive/2026/texmf-dist/tex/latex/l3backend/l3backend-luatex.def
\end{minted}

\bigskip

Como estamos usando {\lmr Lua\LaTeX}, indica que {\lmr\LaTeX} encontró e inicio el procesamiento de las definiciones del motor de {\lmr\LaTeX\kern.15em3} para {\lmr Lua\LaTeX}.
\end{frame}





\begin{frame}[fragile]
\frametitle{El archivo \texttt{.log}}

\begin{minted}[breaklines]{text}
File: l3backend-luatex.def 2026-02-18 L3 backend support: PDF output (LuaTeX)
\l__color_backend_stack_int=\count282
Inserting `l3color' in `luaotfload.parse_color'.)
\end{minted}

\bigskip

Como estamos usando {\lmr Lua\LaTeX}, muestran registros de la carga de las definiciones del motor de {\lmr\LaTeX\kern.15em3} para {\lmr Lua\LaTeX}. Además del fin del procesamiento del mismo.
\end{frame}






\begin{frame}[fragile]
\frametitle{El archivo \texttt{.log}}

\begin{minted}[breaklines]{text}
(./documento.aux)
\openout1 = documento.aux
\end{minted}

\bigskip

Indica lectura del archivo auxiliar y marcado del archivo auxiliar.
\end{frame}





\begin{frame}[fragile]
\frametitle{El archivo \texttt{.log}}

\begin{minted}[breaklines]{text}
LaTeX Font Info:    Checking defaults for OML/cmm/m/it on input line 7.
LaTeX Font Info:    ... okay on input line 7.
...
LaTeX Font Info:    External font `cmex10' loaded for size
(Font)              <6> on input line 9.
\end{minted}

\bigskip

Indica la búsqueda y carga de las fuentes de {\lmr\LaTeX}.
\end{frame}





\begin{frame}[fragile]
\frametitle{El archivo \texttt{.log}}

\begin{minted}[breaklines,breakanywhere]{text}
[1

{/usr/local/texlive/2026/texmf-var/fonts/map/pdftex/updmap/pdftex.map}]
\end{minted}

\bigskip

Indica la escritura de la primera página del documento. Con configuración adicional cargada durante el procesamiento de la misma.
\end{frame}





\begin{frame}[fragile]
\frametitle{El archivo \texttt{.log}}

\begin{minted}[breaklines]{text}
(./documento.aux)
\end{minted}

\bigskip

Indica lectura del archivo auxiliar.
\end{frame}





\begin{frame}[fragile]
\frametitle{El archivo \texttt{.log}}

\begin{minted}[breaklines]{text}
 ***********
LaTeX2e <2025-11-01>
L3 programming layer <2026-04-28>
 ***********
)
\end{minted}

\bigskip

Indica que se terminó de procesar el documento y se descargaron los motores de memoria.
\end{frame}





\begin{frame}[fragile]
\frametitle{El archivo \texttt{.log}}
{
\scriptsize
\begin{minted}[breaklines]{text}
Here is how much of LuaTeX's memory you used:
 462 strings out of 475892
 100000,460012 words of node,token memory allocated
 406 words of node memory still in use:
   3 hlist, 1 vlist, 1 rule, 2 glue, 3 kern, 1 glyph, 4 attribute, 48 glue_spec,
 4 attribute_list, 1 write nodes
   avail lists: 2:80,3:14,4:4,5:28,6:6,7:218,8:1,9:46,10:3,11:4
 23177 multiletter control sequences out of 65536+600000
 37 fonts using 2300247 bytes
 35i,6n,37p,124b,207s stack positions out of 10000i,1000n,20000p,200000b,200000s
\end{minted}
}

\bigskip

Indica cuántos recursos estaban disponibles y cuántos se utilizaron durante el procesamiento del documento.
\end{frame}





\begin{frame}[fragile]
\frametitle{El archivo \texttt{.log}}
\begin{minted}[breaklines,breakanywhere]{text}
</usr/local/texlive/2026/texmf-dist/fonts/opentype/public/lm/lmroman8-regular.otf>
</usr/local/texlive/2026/texmf-dist/fonts/opentype/public/lm/lmroman12-regular.otf>
</usr/local/texlive/2026/texmf-dist/fonts/opentype/public/lm/lmroman17-regular.otf>
\end{minted}

\bigskip

Indica cuáles son las fuentes que se incrustaron en el archivo PDF.
\end{frame}





\begin{frame}[fragile]
\frametitle{El archivo \texttt{.log}}
\begin{minted}[breaklines]{text}
Output written on documento.pdf (1 page, 8615 bytes).
\end{minted}

\bigskip

Indica cuál es el nombre del archivo de salida, su tamaño en páginas y en disco.
\end{frame}





\begin{frame}[fragile]
\frametitle{El archivo \texttt{.log}}
{
\scriptsize
\begin{minted}[breaklines]{text}
PDF statistics: 29 PDF objects out of 1000 (max. 8388607)
 16 compressed objects within 1 object stream
 0 named destinations out of 1000 (max. 1048576)
 1 words of extra memory for PDF output out of 10000 (max. 100000000)
\end{minted}
}

\bigskip

Indica estadísticas de generación del PDF.
\end{frame}





\section{Errores comunes}



\begin{frame}[fragile]
\frametitle{Comando desconocido}

\begin{minted}[linenos,firstnumber=8]{latex}
\maketitl
¡Hola, Mundo {\LaTeX}!
\end{minted}

\bigskip\pause

\begin{minted}[breaklines]{text}
! Undefined control sequence.
l.8 \maketitl

The control sequence at the end of the top line
of your error message was never \def'ed. If you have
misspelled it (e.g., `\hobx'), type `I' and the correct
spelling (e.g., `I\hbox'). Otherwise just continue,
and I'll forget about whatever was undefined.
\end{minted}
\end{frame}






\begin{frame}[fragile]
\frametitle{Comando desconocido}

\begin{minted}[linenos,firstnumber=8]{latex}
\maketitle
¡Hola, Mundo {\LaTeXq}!
\end{minted}

\bigskip\pause

\begin{minted}[breaklines]{text}
! Undefined control sequence.
l.9 ¡Hola, Mundo {\LaTeXq
                        }!
The control sequence at the end of the top line
of your error message was never \def'ed. If you have
misspelled it (e.g., `\hobx'), type `I' and the correct
spelling (e.g., `I\hbox'). Otherwise just continue,
and I'll forget about whatever was undefined.
\end{minted}
\end{frame}





\begin{frame}[fragile]
\frametitle{Referencia desconocida}

\begin{minted}[linenos,firstnumber=8]{latex}
\maketitle
¡Hola, Mundo {\LaTeX}!

\ref{eq:missing}
\end{minted}

\bigskip\pause

\begin{minted}[breaklines]{text}
LaTeX Warning: Reference `eq:missing' on page 1 undefined on input line 11.
\end{minted}
\pause
\begin{minted}[breaklines]{text}
LaTeX Warning: There were undefined references.
\end{minted}
\end{frame}





\begin{frame}[fragile]
\frametitle{Fin de un grupo prematuramente}

\begin{minted}[linenos,firstnumber=8]{latex}
\maketitle
¡Hola, Mundo {\LaTeX!
\end{minted}

\bigskip\pause

\begin{minted}[breaklines]{text}
(\end occurred inside a group at level 1)

### simple group (level 1) entered at line 9 ({)
### bottom level
\end{minted}
\end{frame}




\begin{frame}[fragile]
\frametitle{Entorno sin fin}

\begin{minted}[linenos,firstnumber=8]{latex}
\maketitle

\begin{center}
    ¡Hola, Mundo {\LaTeX}!
\end{minted}

\bigskip\pause

{
\scriptsize
\begin{minted}[breaklines]{text}
! LaTeX Error: \begin{center} on input line 10 ended by \end{document}.

See the LaTeX manual or LaTeX Companion for explanation.
Type  H <return>  for immediate help.
 ...

l.12 \end{document}

Your command was ignored.
Type  I <command> <return>  to replace it with another command,
or  <return>  to continue without it.
\end{minted}
}
\end{frame}





\begin{frame}[fragile]
\frametitle{Entorno desconocido}

\begin{minted}[linenos,firstnumber=8]{latex}
\maketitle

\begin{nonexist}
    ¡Hola, Mundo {\LaTeX}!
\end{nonexists}
\end{minted}

\bigskip\pause

\begin{onlyenv}<1-2>
    \scriptsize
    \begin{minted}[breaklines]{text}
! LaTeX Error: Environment nonexist undefined.

See the LaTeX manual or LaTeX Companion for explanation.
Type  H <return>  for immediate help.
 ...

l.10 \begin{nonexist}

Your command was ignored.
Type  I <command> <return>  to replace it with another command,
or  <return>  to continue without it.
    \end{minted}
\end{onlyenv}

\begin{onlyenv}<3>
    \scriptsize
    \begin{minted}[breaklines]{text}
! LaTeX Error: \begin{document} ended by \end{nonexist}.

See the LaTeX manual or LaTeX Companion for explanation.
Type  H <return>  for immediate help.
 ...

l.12 \end{nonexist}

Your command was ignored.
Type  I <command> <return>  to replace it with another command,
or  <return>  to continue without it.
    \end{minted}
\end{onlyenv}
\end{frame}




\begin{frame}[fragile]
\frametitle{Grupo no cerrado}

\begin{minted}[linenos,firstnumber=8]{latex}
\maketitle

\begin{center}
    ¡Hola, Mundo {\LaTeX!
\end{center}
\end{minted}

\bigskip\pause

{
\scriptsize
\begin{minted}[breaklines]{text}
! Missing } inserted.
<inserted text>
}
l.12 \end{center}

I've inserted something that you may have forgotten.
(See the <inserted text> above.)
With luck, this will get me unwedged. But if you
really didn't forget anything, try typing `2' now; then
my insertion and my current dilemma will both disappear.
\end{minted}
}
\end{frame}





\begin{frame}[fragile]
\frametitle{Advertecia por Espacio}

\begin{minted}[linenos,firstnumber=8]{latex}
\maketitle
¡Hola, Mundo {\LaTeX}! \\

Este es un nuevo párrafo.
\end{minted}

\bigskip\pause

\begin{minted}[breaklines]{text}
Underfull \hbox (badness 10000) in paragraph at lines 9--10

 []
\end{minted}
\end{frame}





\begin{frame}
\frametitle{Estrategias para depuración}

\begin{itemize}[<+->]
    \item Leer los registros (\texttt{.log}).
    \item Comentar código (\texttt{\%}). Comentar las líneas con el problema, analizar si cambia el comportamiento de {\lmr\LaTeX}.
    \item Generar un MWE. Intentar aislar el problema, ver si es posible reproducirlo en un documento más pequeño.
    \item Preguntar. Si no pueden solucionarlo por ustedes mismos, en la web, o ahora con los LLM les pueden ayudar.
\end{itemize}
\end{frame}





\begin{frame}[allowframebreaks]
\frametitle{Referencias}
\nocite{*}
\printbibliography
\end{frame}
\end{document}
